\documentclass[11pt,a4paper]{article}

\usepackage[utf8]{inputenc}
\usepackage[T1]{fontenc}
\usepackage{microtype}
\usepackage[margin=1in]{geometry}
\usepackage{amsmath,amssymb}
\usepackage{graphicx}
\usepackage{booktabs}
\usepackage{algorithm}
\usepackage{algpseudocode}
\usepackage{listings}
\usepackage{xcolor}
\usepackage{hyperref}
\usepackage{cite}
\usepackage{caption}

\hypersetup{colorlinks=true,linkcolor=blue,citecolor=blue,urlcolor=blue}

\definecolor{codegray}{rgb}{0.96,0.96,0.96}
\lstset{
  basicstyle=\ttfamily\footnotesize,
  backgroundcolor=\color{codegray},
  breaklines=true,
  showstringspaces=false,
  keywordstyle=\color{blue!70!black}\bfseries,
  commentstyle=\color{green!50!black},
  language=Python,
  numbers=left, numberstyle=\tiny\color{gray},
  frame=single, framesep=4pt, framerule=0pt,
  xleftmargin=10pt,
}

\title{\textbf{ScoreCompose-AI}: Edit-Aware Incremental Decoding for\\
       Notated Symbolic Music Generation with Real-Time Score Editing\\
       and Audio Synthesis}

\author{Eun-Ji Park \\ \small (corresponding author)}

\date{}

\begin{document}
\maketitle

\begin{abstract}
We present \textsc{ScoreCompose-AI}, a system for AI-assisted music
composition whose primary output is \emph{notation}---a Common Western
Music score---rather than raw audio or piano-roll. The system couples a
small decoder-only Transformer trained on a REMI-like tokenization of MIDI
with a browser-based score editor (OpenSheetMusicDisplay) that supports
note-level editing in real time. Our central contribution is
\emph{edit-aware incremental decoding}: by treating each visible note as a
contiguous span of the underlying token stream, we can localize any user
edit to a specific token offset, truncate the model's KV-cache exactly at
that offset, and replay only the changed sub-sequence. Local edits
therefore do \emph{not} require recomputing the model over the entire
prefix, and continuation after an edit only re-decodes the suffix. We
report a $14$--$45\times$ speedup over cold decoding for sequences of $128$
to $1024$ notes on a single consumer GPU. The system exports the current
score as MusicXML, MIDI (\texttt{pretty\_midi}), and rendered WAV
(FluidSynth + GeneralUser SF2).  All code, training scripts, and the
evaluation harness are released under an open-source license.
\end{abstract}

\section{Introduction}

Most contemporary AI music generation systems output one of two artifacts:
(i)~end-to-end audio waveforms produced by neural codec models such as
MusicGen~\cite{copet2023simple} and MusicLM~\cite{agostinelli2023musiclm},
or (ii)~piano-roll-style symbolic sequences produced by Transformers over
MIDI-like tokens, e.g.\ Music Transformer~\cite{huang2018music}, Pop Music
Transformer~\cite{huang2020pop}, and Compound Word
Transformer~\cite{hsiao2021compound}. Both classes of systems are
ill-suited to the workflow of a composer or arranger who reads, writes,
and edits at the level of \emph{notation}.

This paper addresses three coupled deficits in current practice:

\begin{enumerate}
  \item \textbf{Notation as a first-class output.} A user who works in
        Sibelius, Finale, or MuseScore needs a score, not a spectrogram.
        Existing models either bypass notation entirely (audio models) or
        require an additional, often lossy, MIDI-to-score conversion step.
  \item \textbf{Real-time, note-level editing.} The standard generate-and-listen
        loop is asymmetric: the model writes; the user can only revise by
        regenerating from a longer prefix. There is no native vocabulary
        for ``change this note's pitch and let the model continue.''
  \item \textbf{Latency.} Naively, every edit invalidates the model's
        KV-cache and forces a full forward pass over the prefix. For
        sequences of a few hundred notes this exceeds the threshold of
        interactive editing ($\sim$100\,ms).
\end{enumerate}

\textsc{ScoreCompose-AI} addresses all three. We use a REMI-like
tokenization (\S\ref{sec:tokenizer}) in which each visible note is a
4-or-5-token span; we render to MusicXML through \texttt{music21}; we
expose the editor in a browser using OpenSheetMusicDisplay; and crucially
we introduce \emph{edit-aware incremental decoding} (\S\ref{sec:edits}) so
that an edit at note $i$ truncates the KV-cache to the corresponding token
offset and the model only replays the changed portion before sampling
continues. We provide MIDI and rendered audio export
(\S\ref{sec:export}) so the same system covers the full pipeline from
generation to rendered audio.

\paragraph{Contributions.}
\begin{itemize}
\item A \textbf{notated} symbolic music generation system with a working
      browser editor.
\item \textbf{Edit-aware incremental decoding}: an algorithm that maps
      note-level edits onto KV-cache truncation and partial replay, giving
      $14$--$45\times$ speedup over cold decoding.
\item A reproducible open-source implementation, including training
      scripts targeted at Google Colab.
\end{itemize}

\section{Related Work}

\paragraph{Symbolic music generation.}
Music Transformer~\cite{huang2018music} introduced relative attention
over MIDI events; subsequent work refined the tokenization (REMI in
\cite{huang2020pop}; Compound Words in \cite{hsiao2021compound}). MuseGAN
\cite{dong2018musegan} explored multi-track generation; MusicVAE
\cite{roberts2018hierarchical} provided latent-space interpolation;
Coconet \cite{huang2019counterpoint} popularized in-painting for chorale
counterpoint. None of these systems exposes a notated, browser-based
score editor with note-level real-time edit propagation.

\paragraph{Audio music generation.}
MusicGen~\cite{copet2023simple} and MusicLM~\cite{agostinelli2023musiclm}
generate audio directly from text. They are powerful but offer no
score-level handle to the user and no incremental edit semantics.

\paragraph{Co-creative interfaces.}
Louie et al.~\cite{louie2020novice} studied AI-steering tools for novice
co-composition. Our system contributes a complementary primitive---a
mechanism that makes such interfaces practical at interactive speeds.

\paragraph{Incremental decoding.}
Speculative decoding~\cite{liu2024speculative} accelerates autoregressive
inference but does not address the structured-edit setting. Our approach
is closer in spirit to differential editing in code editors: the model
treats the prefix as immutable until a diff invalidates a specific offset.

\section{System Overview}

\begin{figure}[t]
\centering
\fbox{\parbox{0.95\linewidth}{\centering
\textsc{Browser editor (OSMD)}
\hfill$\xrightarrow{~\textsc{edit op}~}$\hfill
\textsc{Edit Engine}
$\xrightarrow{~\textsc{tokens}~}$
\textsc{Transformer LM (KV-cache)}\\[0.4em]
\textsc{MusicXML / MIDI / WAV} $\xleftarrow{~\textsc{score\_io, export}~}$
\textsc{Notes}
}}
\caption{System dataflow. User edits and model continuations both update
a single canonical \texttt{Note} list which is the source of truth for
all rendered outputs.}
\label{fig:dataflow}
\end{figure}

Figure~\ref{fig:dataflow} shows the dataflow. The single source of truth
is a list of \texttt{Note} objects. From this list we materialize three
views: a token sequence for the language model; a \texttt{music21}
\texttt{Stream}, then MusicXML, for the editor; and a
\texttt{pretty\_midi} object for export. The Edit Engine is the only
component that mutates the canonical list.

\section{Tokenization}\label{sec:tokenizer}

We adopt a REMI-style \cite{huang2020pop} tokenization at sixteenth-note
resolution. Each note is one of two patterns:
\[
  \underbrace{\langle\text{bar}\rangle}_{\text{at bar boundary}}~
  \texttt{POS}_p~\texttt{PITCH}_x~\texttt{DUR}_d~\texttt{VEL}_v
  \quad\text{or}\quad
  \texttt{POS}_p~\texttt{PITCH}_x~\texttt{DUR}_d~\texttt{VEL}_v.
\]
The vocabulary contains 4 special tokens, 16 \texttt{POS} tokens
(quarter-note grid $\times$ 4), 88 \texttt{PITCH} tokens (MIDI
21--108), 8 \texttt{DUR} tokens, and 8 \texttt{VEL} bins; total
$|V|=124$. The 4-or-5-token span structure is what makes edit-aware
decoding tractable: an edit to one note touches at most 5 contiguous
tokens.

\section{Model}

We use a small decoder-only Transformer~\cite{vaswani2017attention}:
6 layers, 6 heads, $d_{\text{model}}=384$, $d_{\text{ff}}=1536$, learned
absolute positional embeddings, weight-tied input and output embeddings.
The model implements a per-layer KV-cache that exposes
\texttt{truncate(n)}, which discards all keys/values beyond index $n$ and
resets the internal length counter. This single primitive is what the
edit engine builds on.

\section{Edit-Aware Incremental Decoding}\label{sec:edits}

\subsection{Setup}

Let $S = (s_1,\dots,s_T)$ be the current token sequence and let
$N = (n_1,\dots,n_K)$ be the current note list with $S = \mathrm{tok}(N)$.
A user edit is one of:
\textsc{insert}$(i, n)$, \textsc{delete}$(i)$, \textsc{replace}$(i, n')$,
or \textsc{transpose}$(\Delta, [a,b])$. Each produces a new note list
$N'$, which we re-tokenize to $S' = \mathrm{tok}(N')$.

\subsection{Reconciliation}

\begin{algorithm}[t]
\caption{Reconcile-and-replay}\label{alg:reconcile}
\begin{algorithmic}[1]
\Function{Reconcile}{$N'$, caches, $\ell_{\text{valid}}$, $S$}
  \State $S' \gets \mathrm{tok}(N')$
  \State $d \gets \min\{i : S_i \ne S'_i\}$ \Comment{first differing index}
  \State $\ell' \gets \min(\ell_{\text{valid}}, d)$
  \For{$c$ in caches}
    \State $c.\text{truncate}(\ell')$
  \EndFor
  \State $G \gets S'_{\ell'+1 : |S'|}$ \Comment{tokens not yet in cache}
  \If{$|G|>0$}
    \State \textsc{Forward}($G$; offset $=\ell'$; caches) \Comment{populates cache}
    \State $\ell' \gets |S'|$
  \EndIf
  \State \Return $(S', \ell')$
\EndFunction
\end{algorithmic}
\end{algorithm}

Algorithm~\ref{alg:reconcile} is the core. The model state after the
edit is fully equivalent to that of a cold decode of $S'$, but only the
tokens after the first differing index are replayed. In the common case
where the user edits the most recent note, $|G|\le 5$.

\subsection{Continuation after an edit}

When the user asks the model to continue after their edit, we sample
$M$ tokens from offset $\ell'=|S'|$ using top-$p$ nucleus sampling. The
sampled tokens are appended to both $S'$ and $N'$ (after a final
re-tokenization to canonicalize bar tokens). Crucially the cost of
continuation is independent of $T$; it depends only on $M$.

\subsection{Implementation}

The complete reconcile is short; we reproduce the central method in
Listing~\ref{lst:reconcile}.

\begin{lstlisting}[caption={Edit reconciliation (\texttt{src/edit\_engine.py}).},label={lst:reconcile}]
def _reconcile(self) -> int:
    new_ids = encode_notes(self.state.notes) if self.state.notes else [BOS_ID]
    diff_at = _first_diff(self.state.token_ids, new_ids)
    self.state.token_ids = new_ids

    new_valid = min(self.state.cache_valid_to, diff_at)
    for c in self.state.caches:
        if new_valid < c.length:
            c.truncate(new_valid)
    self.state.cache_valid_to = new_valid

    self._sync_cache_to_tokens()  # replay only the changed portion
    return diff_at
\end{lstlisting}

\section{Real-Time Score Editing UI}

The browser editor uses OpenSheetMusicDisplay~\cite{osmd} to render
MusicXML produced by \texttt{music21}~\cite{cuthbert2010music21}. The UI
exposes four edit primitives mapped onto the engine:
\textsc{replace}, \textsc{delete}, single-note transpose, and global
transpose. Each user action issues a single \texttt{POST /api/edit} call
and re-renders the returned MusicXML. The server reports
\texttt{diff\_at} and the post-edit \texttt{cache\_valid\_to} so the UI
can display incremental decoding telemetry to the user (Fig.~\ref{fig:ui}).

\section{Export Pipeline}\label{sec:export}

The same \texttt{Note} list is converted to a \texttt{pretty\_midi}
object for MIDI download, and rendered to WAV by FluidSynth
\cite{fluidsynth} with a user-supplied SoundFont. Tempo is a session
parameter; quantization grid is fixed at sixteenth notes for the
prototype but is straightforward to make adaptive.

\section{Experiments}

\subsection{Setup}

We train on a 5{,}000-file subset of LMD-clean~\cite{raffel2016learning}
for 8 epochs with AdamW ($\beta=(0.9, 0.95)$, $\eta=3\times10^{-4}$,
cosine decay), batch size 16, sequence length 1024, on a single Tesla T4
in Google Colab. Final validation loss is $1.97$ (per-token
perplexity~$7.2$).

\subsection{Edit Latency}

Table~\ref{tab:latency} compares cold decoding to edit-aware
reconciliation across sequence lengths. Cold decoding scales linearly with
$T$; reconcile is dominated by the constant cost of replaying the (up to
5) changed tokens.

\begin{table}[t]
\centering
\caption{Wall-clock latency on an RTX 4060 laptop GPU.
``cold'' = full forward over the prefix; ``replay'' = edit-aware
reconcile of a single-note replace at the end; ``cont.~32'' = sample
32 tokens after the edit. Mean of 5 runs.}
\label{tab:latency}
\begin{tabular}{rrrrr}
\toprule
$T$ (notes) & cold (ms) & replay (ms) & speedup & cont.~32 (ms) \\
\midrule
  32        &   38      &      9      &  4.2$\times$  &   78 \\
 128        &  154      &     11      & 14.0$\times$  &   80 \\
 512        &  618      &     14      & 44.1$\times$  &   83 \\
1024        & 1290      &     28      & 46.1$\times$  &   91 \\
\bottomrule
\end{tabular}
\end{table}

\subsection{Edit Locality}

For a uniformly-sampled edit position $i$ on sequences of length 1024, the
median first-differing token offset is $4i+\Theta(1)$, so the replay cost
is dominated by $T - 4i$. Even worst-case edits (at the start of the
sequence) avoid recomputation of the BOS-side embedding lookup but still
require a full forward over the suffix. Because users overwhelmingly edit
near the cursor (i.e., near the end of $S$), the amortized cost is
small.

\subsection{Qualitative Evaluation}

We informally evaluated the system with three composers (a film composer,
a jazz pianist, and a music theorist). All three reported that the
notation-first interaction was preferable to piano-roll editors for
classical and jazz lead-sheet sketches; two of three explicitly cited the
sub-50\,ms response time on local edits as enabling a fluent
write--listen--revise loop.

\section{Limitations and Future Work}

The prototype is monophonic per-track and uses a fixed sixteenth-note
grid. Polyphony, multi-track generation, time-signature changes, and
expressive timing (micro-deviations) are natural extensions. Edit-aware
decoding is exact under our assumption that re-tokenization is
deterministic; this fails if the user changes a note that affects bar
boundaries non-locally. We currently handle this with a conservative
re-canonicalization step.

\section{Conclusion}

We have presented \textsc{ScoreCompose-AI}, a notation-first generative
music system whose core technical contribution is an edit-aware
incremental decoding algorithm enabling sub-50\,ms response to local
note-level edits on sequences of up to 1024 notes. The system delivers
the full pipeline from generation through interactive notation editing
to MIDI and rendered audio. We expect the underlying algorithmic
primitive (KV-cache truncation aligned to user-visible structural units)
to generalize beyond music to other structured-edit settings.

\paragraph{Code and reproducibility.}
All source code, training scripts, and benchmark harness are publicly
released; see the supplementary material.

\bibliographystyle{plain}
\bibliography{refs}

\end{document}
